\section*{Abstract}
AI-text detectors are deployed in high-stakes settings, yet common user behavior can change detector outcomes in ways that are still poorly understood. We study a specific interaction pattern: repeated conversational rejection of drafts as ``too AI-sounding'' versus a single one-shot request to sound less AI-like. Using real API generations from \gptmini, we run a controlled paired experiment on 30 items across two source settings (fresh generations from \hcthree and externally sourced AI news text from \ghostbuster). We evaluate all outputs with two independently trained local detectors (word n-gram and stylometric character n-gram logistic models) with holdout AUC 0.9239 and 0.9538. Contrary to our hypothesis, iterative rejection does not outperform one-shot rewriting in this run. Across all detector/source pairs, the paired mean difference (iterative minus one-shot) is positive, indicating higher AI probability for iterative outputs. The largest effect is on the stylometric detector for external text (+0.0660, 95\% CI [0.0241, 0.1088], FDR-adjusted $p=0.0440$). Pass rates are unchanged on the word detector and drop by 13.3 points on the stylometric detector in both source settings. These results show that rejection loops are not a universally stronger evasion strategy and can push text toward more detectable stylistic regularity.
